\documentclass{article}
\usepackage[preprint]{neurips_2026}

\usepackage[utf8]{inputenc}
\usepackage[T1]{fontenc}
\usepackage{hyperref}
\usepackage{url}
\usepackage{booktabs}
\usepackage{amsfonts}
\usepackage{amsmath}
\usepackage{microtype}
\usepackage{xcolor}

% This is an ES 678 course paper, not a conference submission.
\makeatletter
\renewcommand{\@noticestring}{ES 678 Course Project}
\makeatother

\renewcommand{\arraystretch}{0.9}
\setlength{\tabcolsep}{4pt}

\title{When Accuracy Lies: Detecting Data Leakage in Bank Marketing Models}

\author{%
  Anant Khandelwal \\
  Roll Number: 26210012 \\
  ES 678 \\
  Indian Institute of Technology Gandhinagar
}

\begin{document}
\maketitle

\begin{abstract}
Using the UCI Bank Marketing dataset, we audit a well-known leakage source:
\texttt{duration}, a call-length feature that cannot be known before a call
is placed. We build a 14-feature realistic pre-contact set, excluding
\texttt{duration} and other post-contact information, and compare a
majority-class baseline, Logistic Regression, and Random Forest under a
stratified random split and an approximate chronological split, with all
model and threshold selection confined to validation data and the test set
read only once. Re-adding \texttt{duration} as a leakage diagnostic raises
Average Precision by up to +0.30 (validation) and +0.24 (test), confirmed by
a paired bootstrap whose 95\% confidence intervals exclude zero. The
chronological split further reveals substantial temporal distribution
shift, including a markedly smaller leakage effect on its test set than on
its validation set.
\end{abstract}

\section{Problem and Motivation}
The UCI Bank Marketing dataset records outcomes of a Portuguese bank's
telephone-based term-deposit campaigns \citep{moro2014data,moro2012uci}. Each
row is a client-related campaign contact record; the same client may appear
in more than one row if contacted multiple times, and each row's label
(\texttt{y} = yes/no) indicates whether that specific contact resulted in a
subscription. We study the decision a bank faces \emph{immediately before a
scheduled customer contact begins}: predict whether that contact will result
in a subscription, using only information available at that moment. This
matters because the dataset also includes \texttt{duration}, the call
length, which cannot be known before the call happens yet is strongly
associated with the outcome (0 seconds implies \texttt{y} = no). Prior work
recommends discarding \texttt{duration} for realistic modeling
\citep{moro2014data}, but it is still commonly left in, an instance of the
broader problem of target leakage \citep{kaufman2012leakage}. We ask: (1)
how well can a pre-contact-only model predict subscription, and (2) how much
of commonly reported performance on this dataset is a leakage artifact
rather than genuine signal.

\section{Data and Preparation}
We use \texttt{bank-additional-full.csv} \citep{moro2014data,moro2012uci},
41{,}188 rows and 21 raw columns spanning client attributes, prior-campaign
history, current-campaign contact details, and five economic indicators. A
feature-availability audit classified each column by whether its value would
be known immediately before a scheduled contact, excluding current-call
outcome columns (\texttt{duration}, \texttt{campaign}) and the five economic
indicators (unverifiable at record level), and retaining client-profile fields (\texttt{age}, \texttt{job}, \texttt{marital},
\texttt{education}, \texttt{default}, \texttt{housing}, \texttt{loan}),
prior-campaign history (\texttt{pdays}, \texttt{previous}, \texttt{poutcome}),
and the scheduled contact plan (\texttt{contact}, \texttt{month},
\texttt{day\_of\_week}) -- 13 approved raw features, assumed fixed by the
bank's outreach schedule and known before the call. The sentinel
\texttt{pdays} = 999 (``never previously contacted'') is recoded into a
binary \texttt{previously\_contacted} flag (1{,}515 of 41{,}188 clients) and
\texttt{days\_since\_previous\_contact} (\texttt{NaN} where not applicable),
yielding a final \textbf{14-feature} realistic set (4 numerical, 10
categorical). The literal string \texttt{"unknown"} (e.g.\ \texttt{default}:
8{,}597 rows) is kept as its own category rather than imputed, since
nondisclosure may itself be informative. Remaining missing values use median
or most-frequent imputation, followed by one-hot encoding with unseen
categories ignored at inference; every imputer, encoder, and scaler is fit
on the training split only and never refit on validation/test rows.

\section{Methods}
We compare three classifiers spanning a non-learning floor and two standard
supervised baselines. The \textbf{majority-class baseline} predicts the
single most frequent training-set label for every example, regardless of
input features; its Average Precision therefore equals the positive-class
prevalence in the evaluation set -- a concrete, threshold-free floor any
model with genuine discriminative skill must exceed. \textbf{Logistic
Regression} models $P(y=1 \mid \mathbf{x}) = \sigma(\mathbf{w}^\top
\mathbf{x} + b)$, $\sigma(z) = 1/(1+e^{-z})$, fit by minimizing regularized
log loss with numeric inputs standardized ($z=(x-\mu)/\sigma$, estimated on
training only) and $L_2$ regularization at scikit-learn's default strength
\citep{pedregosa2011scikit}. \textbf{Random Forest} \citep{breiman2001random}
averages $B=100$ bootstrap-trained decision trees with random feature
subsampling; no scaling is applied, since tree splits are scale-invariant.
All models use \texttt{random\_state=42} and otherwise default
hyperparameters, for an honest leakage comparison rather than a tuned
leaderboard entry. To measure leakage directly, we additionally fit
a second, otherwise identical, copy of each real model with \texttt{duration}
added back as a 15th feature; these duration-inclusive models are
\textbf{leakage diagnostics only} and are never treated as deployable.

\section{Evaluation Design}
We construct two independent 70:15:15 train/validation/test splits of the
41{,}188 rows. The \textbf{random split} is stratified on \texttt{y}
($\sim$11.3\% subscription rate preserved in each subset). The
\textbf{chronological split} preserves row order and cuts positionally
(first 70\% / next 15\% / last 15\%), giving an approximate forward-in-time
evaluation on later records rather than an i.i.d.\ sample; its subsets are
\emph{not} class-balanced by construction -- a temporal audit (Section 5)
showed the subscription rate drifts substantially across the timeline.
Preprocessing and models are fit on each split's training rows only. Because
the target class is rare ($\sim$11\%), \textbf{Average Precision} (area
under the precision--recall curve) is the primary metric, since accuracy and
even ROC-AUC can be misleadingly optimistic under imbalance
\citep{saito2015precision}; ROC-AUC, precision, recall, and F1 are secondary.
All model and threshold selection -- including the F1-maximizing threshold
-- uses validation data only, via a grid search over thresholds 0.05--0.95;
thresholds are frozen before the test set is read even once.

\section{Results and Error Analysis}
Validation and test results are reported separately; they are not
interchangeable, since class balance differs between them (notably for the
chronological split, Table~3).

\begin{table}[ht]
\centering
\footnotesize
\caption{Validation results, realistic (no-\texttt{duration}) models.}
\begin{tabular}{llccc}
\toprule
Split & Model & AP & ROC-AUC & F1 \\
\midrule
Stratified & Baseline & 0.113 & 0.500 & 0.000 \\
Stratified & Log.\ Reg. & 0.399 & 0.754 & 0.267 \\
Stratified & Rand.\ Forest & 0.336 & 0.737 & 0.328 \\
Chronological & Baseline & 0.107 & 0.500 & 0.000 \\
Chronological & Log.\ Reg. & 0.198 & 0.657 & 0.000 \\
Chronological & Rand.\ Forest & 0.187 & 0.612 & 0.078 \\
\bottomrule
\end{tabular}
\end{table}

Every baseline's AP equals its split's prevalence, confirming the floor
defined in Section 3. Both real models clearly exceed it on AP in both
splits. Logistic Regression attains marginally higher AP than Random
Forest, but at the default 0.5 threshold its F1 is worse or zero --
chronological Logistic Regression predicts the negative class for every row
(F1 = 0), despite AP (0.198) exceeding baseline, a ranking/classification
gap discussed in Section 6.

\begin{table}[ht]
\centering
\footnotesize
\caption{Test results at the frozen validation-selected threshold.}
\begin{tabular}{llcccc}
\toprule
Split & Model & Prev. & AP & ROC-AUC & F1 \\
\midrule
Strat. & Baseline & 0.113 & 0.113 & 0.500 & 0.000 \\
Strat. & Log.\ Reg. & 0.113 & 0.435 & 0.783 & 0.465 \\
Strat. & Rand.\ Forest & 0.113 & 0.350 & 0.740 & 0.394 \\
Chron. & Baseline & 0.385 & 0.385 & 0.500 & 0.000 \\
Chron. & Log.\ Reg. & 0.385 & 0.468 & 0.605 & 0.416 \\
Chron. & Rand.\ Forest & 0.385 & 0.462 & 0.605 & 0.489 \\
\bottomrule
\end{tabular}
\end{table}

The chronological test prevalence (38.5\%) is far higher than its own
validation prevalence (10.7\%): a temporal audit of 10 equal row-index
segments showed subscription rate rising from 2.8\% to 45.9\%, with the
largest jump (+30.2 points) falling inside the test portion. The stratified
split shows no such shift. This is a property of row ordering, not a claim
about the drift's cause.

\textbf{Duration leakage.} Adding \texttt{duration} (diagnostics only)
produces large, consistent gains on both validation and test:

\begin{table}[ht]
\centering
\footnotesize
\caption{$\Delta$AP / $\Delta$ROC-AUC from adding \texttt{duration}.}
\begin{tabular}{llcccc}
\toprule
& & \multicolumn{2}{c}{Validation} & \multicolumn{2}{c}{Test} \\
Split & Model & $\Delta$AP & $\Delta$AUC & $\Delta$AP & $\Delta$AUC \\
\midrule
Strat. & Log.\ Reg. & +0.174 & +0.161 & +0.128 & +0.136 \\
Strat. & Rand.\ Forest & +0.269 & +0.182 & +0.236 & +0.174 \\
Chron. & Log.\ Reg. & +0.231 & +0.220 & +0.078 & +0.097 \\
Chron. & Rand.\ Forest & +0.300 & +0.270 & +0.068 & +0.081 \\
\bottomrule
\end{tabular}
\end{table}

\textbf{Failure modes (validation, realistic models).} False-negative rates
are high under class imbalance (74--83\% for the stratified split);
chronological Logistic Regression has zero true positives at the default
threshold. The most notable subgroup finding is \texttt{poutcome=success}:
over-predicted positive on the stratified split (Random Forest
false-positive rate 60\%), but almost entirely missed on the chronological
split (Section 6, with counts).

\section{Graduate Depth Analysis}
\textbf{Paired bootstrap on validation.} For each split and model, we
resampled the matched validation predictions (with- and without-
\texttt{duration}, same rows) 1{,}000 times and recomputed $\Delta$AP on
each resample, yielding a 95\% percentile confidence interval.

\begin{table}[ht]
\centering
\footnotesize
\caption{Paired bootstrap $\Delta$AP, 95\% CI, validation only (n=6{,}178).}
\begin{tabular}{llcc}
\toprule
Split & Model & Observed $\Delta$AP & 95\% CI \\
\midrule
Strat. & Log.\ Reg. & +0.174 & [+0.134, +0.210] \\
Strat. & Rand.\ Forest & +0.269 & [+0.236, +0.299] \\
Chron. & Log.\ Reg. & +0.231 & [+0.196, +0.269] \\
Chron. & Rand.\ Forest & +0.300 & [+0.265, +0.341] \\
\bottomrule
\end{tabular}
\end{table}

All four intervals exclude zero: the AP increase from \texttt{duration} is
not attributable to resampling noise on validation data. This interval
describes only validation-set sampling uncertainty; no equivalent bootstrap
was computed on the test set, and it should not be read as a statement
about test-set uncertainty.

\textbf{Threshold sensitivity: ranking vs.\ classification.} AP and ROC-AUC
are threshold-independent ranking metrics; precision, recall, and F1 require
a fixed decision threshold. Sweeping validation thresholds from 0.05 to 0.95
shows the default 0.5 is rarely optimal: chronological Random Forest's
validation F1 rises from 0.078 (at 0.5) to 0.251 (at 0.15); chronological
Logistic Regression rises from 0.000 to 0.274 (at 0.10). A model can
therefore rank positives above negatives reasonably well while appearing to
have ``learned nothing'' under an unexamined default threshold.

\textbf{Chronological leakage: validation vs.\ test.} The leakage effect
(Table 3) shrinks far more on the chronological split than the stratified
split from validation to test: Random Forest $\Delta$AP falls from +0.300 to
+0.068, a $\sim$77\% relative drop, versus $\sim$12\% for the stratified
split. This coincides with the chronological test set's much higher
prevalence (Section 5) and indicates that, for this split, validation-set
leakage magnitude is not a reliable estimate of test-set leakage magnitude
-- an association reported without attributing a causal mechanism.

\textbf{Subgroup failure findings (validation, with counts).} For
\texttt{poutcome=success}: stratified Random Forest has a 60.0\%
false-positive rate ($n{=}187$; 65 actual negatives, 39 false positives),
while chronological Random Forest has a 100\% false-negative rate on the
same category ($n{=}67$ actual positives, 67 missed, only 1 false positive).
By \texttt{job} (subgroups $\ge$100 rows), stratified Random Forest's
highest false-positive rate is in \texttt{student} (12.2\%, $n{=}132$) and
\texttt{retired} (11.3\%, $n{=}237$); under the chronological split, every
\texttt{job} subgroup shows a 100\% false-negative, 0\% false-positive rate
for Logistic Regression, consistent with its all-negative predictions.

\section{Limitations and Conclusion}
This project's central finding is that duration-inclusive models
substantially overstate performance relative to a realistic pre-contact
model, on both validation and the held-out test set, confirming this
dataset's well-documented prediction-time leakage. The chronological split
further shows that performance and leakage do not generalize uniformly over
time, since its test set's much higher subscription rate shrinks the
measured leakage relative to validation and shifts the operating point. The
dataset provides no complete timestamp or customer identifier, so the
chronological split uses row order as an \emph{approximate} proxy for time
rather than verified dates, and apparent temporal drift may partly reflect
unobserved record-ordering effects. The scheduled-contact features
(\texttt{contact}, \texttt{month}, \texttt{day\_of\_week}) are assumed known
before the call, assuming they reflect bank-side scheduling, which the data
cannot verify. The reported 95\% confidence intervals come from a paired
bootstrap on \emph{validation data only}, with no equivalent interval for
test-set leakage, reported only as a point estimate. Finally, every result
here is a predictive association, not a causal claim about why any client
subscribes or about the temporal shift's cause. Within these bounds, a
meaningful share of commonly reported bank-marketing performance is a
leakage artifact, not genuine pre-contact predictive skill.

\section{Reproducibility and AI-Use Disclosure}
\textbf{Data.} \texttt{bank-additional-full.csv}, UCI Machine Learning
Repository, \url{https://doi.org/10.24432/C5K306} (41{,}188 rows). All
experiments use \texttt{random\_state=42}. Two fixed 70:15:15
train/validation/test splits (stratified random; chronological by row
order) are saved as index files and reused unchanged. Every preprocessing
step is fit on the training split only, and classification thresholds are
selected on validation data before the test set is read even once. The full
pipeline -- data audit, preprocessing, model training, leakage diagnostics,
bootstrap analysis, and final test evaluation -- is implemented across nine
Jupyter notebooks (\texttt{01\_data\_understanding} through
\texttt{09\_final\_test\_evaluation}) in the project repository.

\textbf{AI-use disclosure.} ChatGPT and Claude assisted with: (1)
identifying and selecting a suitable publicly available dataset; (2)
step-by-step project planning and workflow design; (3) writing and
implementing code; (4) identifying and debugging coding errors; (5)
double-checking implementation, experimental procedures, and results; and
(6) understanding machine-learning concepts and drafting this paper's
documentation. The author independently reviewed all code, executed every
experiment, examined the resulting outputs, and takes full responsibility
for the final work.

\begin{thebibliography}{9}

\bibitem[Moro et al.(2014)]{moro2014data}
S.~Moro, P.~Cortez, and P.~Rita.
\newblock A data-driven approach to predict the success of bank telemarketing.
\newblock \emph{Decision Support Systems}, 62:22--31, 2014.

\bibitem[Moro et al.(2012)]{moro2012uci}
S.~Moro, P.~Rita, and P.~Cortez.
\newblock Bank Marketing [Dataset].
\newblock \emph{UCI Machine Learning Repository}, 2012.
\newblock \url{https://doi.org/10.24432/C5K306}.

\bibitem[Kaufman et al.(2012)]{kaufman2012leakage}
S.~Kaufman, S.~Rosset, C.~Perlich, and O.~Stitelman.
\newblock Leakage in data mining: Formulation, detection, and avoidance.
\newblock \emph{ACM Transactions on Knowledge Discovery from Data}, 6(4):1--21, 2012.

\bibitem[Breiman(2001)]{breiman2001random}
L.~Breiman.
\newblock Random forests.
\newblock \emph{Machine Learning}, 45(1):5--32, 2001.

\bibitem[Pedregosa et al.(2011)]{pedregosa2011scikit}
F.~Pedregosa, G.~Varoquaux, A.~Gramfort, et al.
\newblock Scikit-learn: Machine learning in Python.
\newblock \emph{Journal of Machine Learning Research}, 12:2825--2830, 2011.

\bibitem[Saito and Rehmsmeier(2015)]{saito2015precision}
T.~Saito and M.~Rehmsmeier.
\newblock The precision-recall plot is more informative than the ROC plot
  when evaluating binary classifiers on imbalanced datasets.
\newblock \emph{PLOS ONE}, 10(3):e0118432, 2015.

\end{thebibliography}

\end{document}
